% Vorlage für WABs der Provadis Hochschule
% basiert auf "Universität Ulm Praktikumsbericht Vorlage"
% von Max Sch.
% CC BY 4.0
%
% Berichtsinfos.tex
%
% Hier werden individuelle Inforamtionen festgelegt.

% Art der Arbeit (genau eine Zeile aktiv lassen)
% Pflichtangabe Deckblatt, siehe Anforderungen "Formatierung 1a"
\newcommand{\artderarbeit}{Wissenschaftlich angeleitete Berufspraxis}
%\newcommand{\artderarbeit}{Wissenschaftlicher Kurzbericht}
%\newcommand{\artderarbeit}{Bachelorarbeit}
%\newcommand{\artderarbeit}{Masterarbeit}

% Nur bei Praxisberichten (WAB / Kurzbericht): das Fach. Bei BT/MT leer lassen: {}
\newcommand{\fach}{Name des Fachs / der Veranstaltung}

% Nur bei BT/MT: angestrebter akademischer Grad. Bei Praxisberichten leer lassen: {}
\newcommand{\abschluss}{}
%\newcommand{\abschluss}{Bachelor of Science (B.Sc.)}
%\newcommand{\abschluss}{Master of Science (M.Sc.)}

%persönliche Daten
\newcommand{\fullname}{First Last}
\newcommand{\email}{first.last@stud-provadis-hochschule.de}
\newcommand{\matnr}{XXXX}
\newcommand{\studiengang}{Business Information Management}
\newcommand{\studiengruppe}{BBIM26}
% Telefonnummer ist optional. Zum Weglassen leer lassen: {}
\newcommand{\telMobil}{+49 XXX XXXXXXX}
\newcommand{\fakultaet}{Informatik und Wirtschaftsinformatik}

%Daten zum Unternehmen
\newcommand{\titel}{Textvorlage für wissenschaftliche Arbeiten\\
Titel und Untertitel der Arbeit}
\newcommand{\unternehmen}{Musterfirma}

\newcommand{\abgabedatum}{XX.XX.XXXX}
\newcommand{\ort}{Frankfurt am Main}

% Optional: Bild der digitalen Unterschrift für die Selbstständigkeitserklärung,
% z.B. {Bilder/Unterschrift.png}. Bei Upload im HVS optional -> leer lassen: {}
\newcommand{\unterschrift}{}

% Daten zu Betreuern
% Referent:in = betreuende Dozent:in
\newcommand{\refFirst}{Prof. Dr. First Last}
% Korreferent:in nur bei Bachelor- und Masterarbeiten, sonst leer lassen: {}
\newcommand{\refSecond}{Prof. Dr. First Last}

% Zitierweise: Kurzzitierweise nach Harvard in Fußnoten (typisch in der Wirtschaftsinformatik,
% "siehe Weber (2020): S. 8.") oder numerisch (typisch in der Kerninformatik, IEEE-ähnlich,
% "[1, S. 8]").
\newif\ifharvard
\harvardtrue
%\harvardfalse

% Wurde KI verwendet? Steuert den Inhalt des (verpflichtenden) KI-Verzeichnisses.
\newif\ifKIverwendet
\KIverwendettrue
%\KIverwendetfalse


%Deklariert den Namen des Authors, sowie den Namen der generierten pdf Datei
\title{\artderarbeit\ -- \unternehmen}
\author{\fullname}
\date{\abgabedatum}
